% !TeX program = pdflatex
% !BIB program = bibtex
\PassOptionsToPackage{spanish}{babel}
\documentclass[final,1p,times,authoryear]{elsarticle}
\makeatletter
\def\ps@pprintTitle{%
  \let\@oddhead\@empty
  \let\@evenhead\@empty
  \let\@oddfoot\@empty
  \let\@evenfoot\@empty
}
\makeatother
\usepackage[utf8]{inputenc}
\usepackage[T1]{fontenc}
\usepackage{amsmath}
\usepackage{amsthm}
\usepackage{amssymb}
\usepackage{newtxtext,newtxmath}
\usepackage[spanish]{babel}
\usepackage[babel=true]{csquotes}
\usepackage{tcolorbox}
\usepackage{booktabs}
\usepackage{multirow}
\usepackage{float}
\usepackage[unicode=true,pdfencoding=auto]{hyperref}
\usepackage{threeparttable}
\usepackage{tikz}
\usepackage{array, longtable, tabularx}
\usepackage{adjustbox}
\usepackage{ltxtable}
\usepackage{makecell}
\usepackage{pdflscape}
\usepackage{enumitem}
\usepackage{wrapfig}
\usepackage[expansion=false]{microtype}
\usepackage{mathtools}
\usepackage{bookmark}
\usepackage[spanish]{cleveref}

\allowdisplaybreaks
\tolerance=1000
\emergencystretch=10pt
\hyphenpenalty=50
\exhyphenpenalty=50

\setlength{\textfloatsep}{10pt plus 2pt minus 2pt}
\setlength{\floatsep}{8pt plus 2pt minus 2pt}
\setlength{\intextsep}{8pt plus 2pt minus 2pt}
\setlength{\abovecaptionskip}{6pt}
\setlength{\belowcaptionskip}{6pt}
\setlength{\columnsep}{20pt}
\setlength{\parskip}{6pt plus 1pt minus 1pt}
\setlength{\parindent}{0pt}

\providecommand{\X}{X}
\providecommand{\T}{T}
\providecommand{\DD}{\mathrm{DD}}
\providecommand{\DI}{\mathrm{DI}}
\providecommand{\DT}{\mathrm{DT}}
\providecommand{\Lmax}{L_{\max}}
\providecommand{\epsedge}{\epsilon_{\text{edge}}}
\providecommand{\epscontrib}{\epsilon_{\text{contrib}}}

\begin{document}

\begin{frontmatter}
\title{Metodología del Índice de Dependencia Económica Elcano (IDEE): Redes Críticas, Dependencias Indirectas y Poder Estructural en el Comercio Global}

\author[upo,elcano]{Manuel Hidalgo-Pérez\corref{cor1}}\ead{mhidalgo@rielcano.org}
\cortext[cor1]{Autor para correspondencia: Manuel Hidalgo-Pérez, mhidalgo@rielcano.org}

\affiliation[upo]{
  organization={Universidad Pablo de Olavide},
  addressline={Ctra Utrera s/n},
  postcode={41013},
  city={Sevilla},
  country={España}
}

\affiliation[elcano]{
  organization={Fundación Real Instituto Elcano},
  addressline={C. de Príncipe de Vergara, 51},
  postcode={28006},
  city={Madrid},
  country={España}
}

\begin{abstract}
La creciente fragmentación geopolítica ha transformado el panorama del comercio internacional, exponiendo vulnerabilidades económicas previamente subestimadas. Este trabajo propone un nuevo indicador de seguridad económica que supera las limitaciones de los enfoques tradicionales al capturar sistemáticamente las dependencias comerciales indirectas transmitidas a través de países intermediarios. Mediante la integración de matrices Input-Output, teoría de redes y algoritmos de propagación de shocks, nuestro indicador cuantifica tanto las dependencias directas como las indirectas, revelando vulnerabilidades ocultas en las cadenas de suministro globales. El estudio empírico demuestra que los análisis tradicionales subestiman significativamente la exposición a disrupciones externas en sectores estratégicos. Estos hallazgos tienen implicaciones importantes para las políticas de autonomía estratégica, sugiriendo que una diversificación efectiva debe considerar toda la estructura de la red comercial global en lugar de limitarse a las relaciones bilaterales directas.
\end{abstract}

\begin{keyword}
Seguridad económica \sep Dependencias comerciales \sep Teoría de redes \sep Autonomía estratégica \sep Cadenas globales de valor\\
\textbf{Códigos JEL:} F14, F15, F52, C67, D85
\end{keyword}
\end{frontmatter}

\section{Introducción}

En un contexto de creciente fragmentación geoeconómica, tensiones sistémicas entre las principales potencias y disrupciones frecuentes en las cadenas de suministro globales, la evaluación de la seguridad económica nacional requiere herramientas analíticas más refinadas. La literatura existente se ha centrado predominantemente en métricas agregadas de comercio bilateral o dependencia directa, sin capturar las interdependencias estructurales que surgen a través de las redes de producción global y el comercio indirecto.

Este documento aborda esta brecha mediante el desarrollo de un Índice de Dependencia Económica Elcano (IDEE), diseñado para medir la exposición de un país a disrupciones exógenas mediante la integración de dependencias directas e indirectas, y considerando la relevancia estructural de los intermediarios críticos.

Para los propósitos de esta investigación, definimos la seguridad económica como "la capacidad de una nación para mantener la integridad de sus funciones económicas esenciales y sus cadenas de suministro críticas frente a disrupciones exógenas, ya sea que estas surjan de desastres naturales, fallos tecnológicos o acciones deliberadas de otros actores".

Esta definición distingue la seguridad económica de conceptos relacionados:
\begin{itemize}
    \item \textbf{Soberanía económica:} A diferencia de la soberanía, que a menudo implica autarquía o una dependencia mínima de insumos extranjeros, la seguridad económica reconoce la realidad y los beneficios de la interdependencia mientras se centra en la gestión de vulnerabilidades críticas.
    \item \textbf{Resiliencia:} Mientras que la resiliencia se refiere a la capacidad de recuperarse de las disrupciones, la seguridad económica abarca tanto la prevención de disrupciones críticas como el mantenimiento de la funcionalidad durante las mismas.
    \item \textbf{Competitividad:} En contraste con la competitividad, que se centra en la ventaja relativa en los mercados globales, la seguridad económica se centra en la estabilidad y fiabilidad de las funciones económicas esenciales para el bienestar nacional.
    \item \textbf{Seguridad nacional tradicional:} A diferencia de las concepciones tradicionales de seguridad nacional, que abordan principalmente amenazas militares, la seguridad económica reconoce que las amenazas existenciales pueden surgir de disrupciones en el suministro, dependencia tecnológica o coerción económica.
\end{itemize}

Nuestro Índice de Dependencia Económica Elcano operacionaliza esta definición midiendo las vulnerabilidades tanto directas como indirectas en todas las industrias, con especial atención a los sectores donde las disrupciones tendrían efectos en cascada en toda la economía. El índice está diseñado para identificar no solo las dependencias obvias en el comercio bilateral directo, sino también las vulnerabilidades ocultas que surgen de la posición de un país dentro de complejas redes de suministro globales.

\section{Fundamentos Teóricos del Índice de Dependencia Económica Elcano}

Este trabajo ancla su enfoque en la medición de la seguridad económica en tres marcos teóricos complementarios: la teoría de la interdependencia compleja, el poder estructural en la economía política internacional y la interdependencia armificada (\textit{weaponized interdependence}).

\textbf{Interdependencia Compleja y Vulnerabilidad de Red}
Basándonos en el trabajo fundamental de \cite{keohane1977power} sobre la interdependencia compleja, conceptualizamos la seguridad económica no simplemente como la ausencia de dependencias bilaterales, sino como una propiedad que surge de la posición de un país dentro de múltiples redes de intercambio superpuestas. Como señalan \cite{farrell2019weaponized}, "la interdependencia ha creado redes complejas con topografías particulares", donde ciertos nodos se convierten en puntos críticos para los flujos globales. Nuestro IDEE operacionaliza este entendimiento basado en redes al modelar explícitamente cómo las disrupciones se propagan a través de múltiples caminos dentro del sistema de comercio global.

\textbf{Poder Estructural en la Economía Política Internacional}
El segundo pilar teórico es el concepto de Strange \citeyear{strange1988states} de poder estructural, entendido como "el poder de dar forma y determinar las estructuras de la economía política global [...] dentro de las cuales otros estados, sus instituciones políticas, sus empresas económicas [...] tienen que operar". En el marco de Strange, el poder estructural opera a través del control sobre la seguridad, la producción, las finanzas y el conocimiento. Nuestro índice operacionaliza específicamente el poder estructural en el dominio de la producción al cuantificar cómo el control sobre nodos críticos en las cadenas de suministro confiere una influencia que se extiende mucho más allá de las relaciones bilaterales directas. Al identificar dependencias indirectas e intermediarios críticos, hacemos visible lo que Strange describió como el "ejercicio indirecto e inobservable del poder" en las relaciones económicas globales.

\textbf{Interdependencia Armificada (\textit{Weaponized Interdependence})}
Finalmente, nuestro trabajo se compromete con la teoría de Farrell y Newman \citeyear{farrell2019weaponized} sobre la "interdependencia armificada", que postula que "los estados con autoridad política sobre los nodos en las estructuras en red internacionales a través de las cuales viajan el dinero, la información y los bienes están posicionados de manera única para imponer costes a otros". Esta teoría distingue entre "efectos panópticos" (la capacidad de monitorear flujos) e "hitos o puntos de estrangulamiento" (el efecto \textit{chokepoint}, la capacidad de restringir flujos). Nuestro IDEE mide específicamente la vulnerabilidad a estos últimos.

\section{Motivación y Contexto Geoeconómico}

La seguridad económica ha surgido como una prioridad estratégica global en un contexto marcado por sucesivas crisis sistémicas que han revelado vulnerabilidades críticas en las cadenas de suministro internacionales. Los últimos años han mostrado cómo una dependencia excesiva de proveedores específicos puede socavar no solo la estabilidad económica sino también la seguridad nacional en un sentido más amplio.

La pandemia de COVID-19 representó el primer shock sistémico que expuso claramente estas fragilidades. Las disrupciones en la producción y distribución de productos esenciales demostraron que la concentración geográfica de las capacidades de producción puede conducir a escaseces críticas con efectos en cascada en múltiples sectores económicos.

La invasión rusa de Ucrania intensificó aún más estas preocupaciones, revelando vulnerabilidades particularmente agudas en los sectores energético y alimentario europeos. Recientemente, las barreras comerciales se han utilizado como herramientas diplomáticas, lo que subraya cómo las preocupaciones de seguridad nacional están cada vez más entrelazadas con la política comercial. Un ejemplo reciente es el anuncio de aranceles recíprocos por parte de la administración estadounidense sobre importaciones de la Unión Europea y China, lo que subraya la vulnerabilidad de las relaciones comerciales a consideraciones no económicas.

En este contexto, el acceso a materias primas críticas, tecnologías emergentes e infraestructuras digitales se ha vuelto central para la planificación estratégica. Entidades como la Unión Europea han desarrollado estrategias de "autonomía estratégica abierta" para equilibrar los beneficios del libre comercio con la necesidad de asegurar el acceso a recursos esenciales.

\section{Revisión de la Literatura}

\subsection{Evolución de la Medición de la Dependencia Comercial}
La medición de la dependencia comercial ha sido una preocupación central en la economía internacional desde mediados del siglo XX. Los primeros enfoques se centraron principalmente en las relaciones bilaterales y la concentración de importaciones. Sin embargo, a medida que la producción global se fragmentaba a través de las fronteras, la insuficiencia de estas medidas tradicionales se hizo evidente. La integración de las cadenas globales de valor (CGV) ha alterado fundamentalmente la naturaleza de las dependencias comerciales.

\subsection{Enfoques Metodológicos para la Medición de Dependencias}
\subsubsection{Índices de Concentración y Medidas Bilaterales}
Una de las herramientas más utilizadas es el Índice de Herfindahl-Hirschman (HHI), que evalúa la concentración geográfica de las importaciones. Instituciones como la Comisión Europea han empleado este índice para identificar dependencias estratégicas, estableciendo umbrales específicos para señalar riesgos para la seguridad económica. No obstante, estos índices presentan limitaciones: son estáticos y no capturan las dependencias indirectas que surgen a través de terceros países.

\subsubsection{Análisis Input-Output y Vínculos Intersectoriales}
El análisis Input-Output (I-O), iniciado por \cite{Leontief1951}, proporciona un marco para comprender la red de dependencias intersectoriales. La matriz inversa de Leontief, $(I - A)^{-1}$, captura los requisitos tanto directos como indirectos de cada sector para producir una unidad de demanda final. Una extensión importante es el concepto de "longitudes de propagación promedio" (APL) desarrollado por \cite{Dietzenbacher2005}, que mide la distancia económica entre sectores a través de vínculos hacia adelante y hacia atrás.

\subsubsection{Análisis de Redes y Medidas de Centralidad}
El análisis de redes emerge como un enfoque complementario que modela el comercio como un sistema complejo de nodos (países o sectores) y enlaces (flujos comerciales) interconectados. Pioneros como \cite{Fagiolo2010} caracterizaron la Red de Comercio Mundial (WTN) documentando sus propiedades estructurales. Conceptos de centralidad como la intermediación (\textit{betweenness}) permiten identificar nodos que actúan como "hubs" estructurales o cuellos de botella críticos.

\subsubsection{Modelos de Propagación de Shocks}
Investigaciones recientes se centran en los mecanismos dinámicos a través de los cuales los shocks económicos se propagan. \cite{Acemoglu2012} y \cite{Carvalho2013} demostraron que los shocks a nivel micro a empresas o sectores individuales pueden tener efectos agregados significativos cuando las redes de producción exhiben ciertas propiedades estructurales, como una distribución de grados con colas pesadas.

\section{Desarrollo del Indicador Propuesto}

\subsection{Conceptualización}
La propuesta de este trabajo es un indicador integral que combina el análisis matricial, la teoría de redes y los algoritmos de propagación para identificar y cuantificar vulnerabilidades ocultas.

\subsubsection{Definición de Dependencia Directa e Indirecta}
La dependencia comercial se define como la exposición de un país a disrupciones de suministro originadas en otro, considerando todos los canales posibles.
\begin{enumerate}
    \item \textbf{Dependencia Directa ($\DD$):} Proporción de importaciones directas en el suministro total de un país.
    \item \textbf{Dependencia Indirecta ($\DI$):} Proporción del suministro atribuible a un exportador a través de rutas comerciales de múltiples pasos (intermediarios).
\end{enumerate}

\subsubsection{Ventajas del Indicador Propuesto}
Nuestro enfoque destaca por tres ventajas clave: (1) identificación de vulnerabilidades sistémicas ocultas que los índices bilaterales no detectan, (2) evaluación cuantitativa de la resiliencia ante disrupciones mediante el modelado de rutas alternativas, y (3) análisis prospectivo de escenarios de fragmentación geoeconómica.

\begin{table}[!t]
\caption{Comparación de Métodos para la Medición de la Dependencia Comercial} \label{tab:comparacion-metodos}
\centering
\footnotesize
\renewcommand{\arraystretch}{1.1}
\begin{adjustbox}{width=\textwidth,center}
    \begin{tabular}{@{}l*{6}{c}@{}}
        \toprule
        \makecell[l]{\textbf{Característica}} & \textbf{HHI} & \makecell{\textbf{Conc.}\\\textbf{Imp.}} & \textbf{TiVA} & \makecell{\textbf{I-O}\\\textbf{Trad.}} & \makecell{\textbf{Re-}\\\textbf{des}} & \makecell{\textbf{Ind.}\\\textbf{Prop.}} \\
        \midrule
        Captura dep. directa & \checkmark & \checkmark & \checkmark & \checkmark & \checkmark & \checkmark \\
        Captura dep. indirecta & \texttimes & \texttimes & Parcial & Solo dirc. & \checkmark & \textbf{Completa} \\
        Normalización (0-1) & \checkmark & \checkmark & \texttimes & \texttimes & \texttimes & \checkmark \\
        Basado en Input-Output & \texttimes & \texttimes & \checkmark & \checkmark & \texttimes & \checkmark \\
        Teoría de redes & \texttimes & \texttimes & \texttimes & \texttimes & \checkmark & \checkmark \\
        Modela propagación shocks & \texttimes & \texttimes & \texttimes & Limitada & Parcial & \checkmark \\
        Detecta vuln. ocultas & \texttimes & \texttimes & Parcial & \texttimes & Parcial & \checkmark \\
        Análisis de resiliencia & \texttimes & \texttimes & Limitado & Limitado & Parcial & \checkmark \\
        \bottomrule
    \end{tabular}
\end{adjustbox}
\end{table}

\subsection{Formulación Matemática}
La dependencia total $\text{DT}_{ij}$ de un importador $j$ respecto a un exportador $i$ se define como la suma de la dependencia directa e indirecta.
Consideramos la matriz de comercio normalizada $\mathbf{T}$, donde cada elemento $T_{ab} = x_{ab} / S_b$ representa la fracción del suministro de $b$ que proviene directamente de $a$.
La fuerza de un camino comercial $p = (i, k_1, \ldots, k_{\ell-1}, j)$ de longitud $\ell$ se calcula como el producto de las dependencias a lo largo de sus aristas:
\begin{equation}
F(p) = \prod_{(a,b) \in \text{edges}(p)} T_{ab} = T_{ik_1} \cdot T_{k_1k_2} \cdots T_{k_{\ell-1}j}
\end{equation}

La dependencia total es la suma de las fuerzas de todos los caminos económicamente significativos de $i$ a $j$:
\begin{equation}
\text{DT}_{ij} = \sum_{\ell=1}^{L_{\max}} \sum_{p \in \mathcal{P}_{ij}^{\ell}} F(p)
\end{equation}

\subsection{Métricas Operativas del Dashboard: Índice de Dependencia e Índice de Vulnerabilidad}

Para alinear el análisis académico con la visualización práctica del portal web de la iniciativa, definimos dos métricas operativas clave en la interfaz de usuario:

\begin{enumerate}
    \item \textbf{Índice de Dependencia (ID):} Este indicador mide la concentración de las importaciones directas de un país en sus socios comerciales. Para facilitar la comprensión del lector general en el portal web (evitando tecnicismos complejos como HHI o "proveedores efectivos"), se construye combinando un \textit{Índice de Herfindahl-Hirschman (HHI) normalizado} con la \textit{cuota de mercado del primer proveedor}. Este índice establece la base de concentración comercial aduanera.
    
    \item \textbf{Índice de Vulnerabilidad (IV):} Para convertir la concentración aduanera directa en una medida de exposición sistémica real, se calcula el \textbf{Índice de Vulnerabilidad}. A nivel sectorial, este coincide con la \textbf{Dependencia Total} ($\text{DT}_{js}$) obtenida mediante el algoritmo de redes (que ya integra de forma natural los caminos directos e indirectos). 
\end{enumerate}

La \textbf{Vulnerabilidad Nacional} agregada ($V_j$) ---que constituye el valor oficial del \textbf{Índice de Dependencia Económica Elcano (IDEE)} a nivel de país--- se obtiene promediando las dependencias totales sectoriales, ponderadas por el peso económico (el valor de las importaciones en dólares) de cada industria sobre el total nacional:
\begin{equation}
V_j = \sum_{s} w_{js} \cdot \text{DT}_{js}
\end{equation}
donde $\text{DT}_{js}$ es la dependencia total del país $j$ en el sector $s$, y $w_{js}$ es la participación del sector $s$ en las importaciones totales de $j$. El índice final se normaliza en un rango de 0 a 100 en la interfaz del portal.

\subsection{Implementación Computacional y Consideraciones Prácticas}

La aplicación práctica a datos reales plantea desafíos computacionales significativos debido al crecimiento exponencial de caminos potenciales. Nuestra implementación incorpora estrategias de optimización:
\begin{itemize}
    \item \textbf{Umbral de relevancia comercial ($\theta$):} Se ignoran los flujos comerciales cuyas contribuciones relativas son inferiores a un umbral $\theta$, lo que poda drásticamente el árbol de búsqueda.
    \item \textbf{Umbral de convergencia ($\varepsilon$):} El cálculo de rutas más largas se detiene cuando su contribución marginal a la dependencia acumulada es inferior a $\varepsilon$.
    \item \textbf{Longitud máxima del camino ($L_{\max}$):} En la práctica, se utilizan valores entre 3 y 5, ya que rutas más largas suelen tener fuerzas insignificantes.
    \item \textbf{Uso de matrices dispersas:} Dado que la mayoría de los países solo comercian con un subconjunto de socios, el uso de estructuras de datos dispersas optimiza drásticamente el uso de memoria y la velocidad de cálculo.
\end{itemize}

\section{Datos y suministro doméstico}\label{sec:datos}

\subsection{Flujos comerciales y denominador del suministro}
El índice se calcula a partir de la \textit{International Trade and Production Database for Estimation}, versión 3 (ITPD-E-R03), publicada por la Comisión de Comercio Internacional de Estados Unidos \citep{Borchert2021}. La base cubre 170 industrias (agricultura, minería y energía, manufacturas y servicios) y 236 economías, y recoge tanto el comercio internacional como el doméstico, definido este como la producción bruta menos las exportaciones. Para cada industria, la matriz de comercio $\mathbf{X}$ contiene por tanto el comercio doméstico en su diagonal ($x_{cc}$) y las importaciones fuera de ella, y el suministro del país $c$ es
\begin{equation}
S_c = x_{cc} + \sum_{a \neq c} x_{ac}.
\end{equation}
El comercio doméstico entra en el índice solo a través de este denominador. Las rutas indirectas excluyen como intermediarios al importador y al exportador, de modo que la diagonal nunca actúa como eslabón de una ruta, pero sí determina qué parte del suministro de cada país depende de fuentes exteriores, tanto para el importador final como para cada intermediario.

\subsection{Lagunas en el comercio doméstico}
El comercio doméstico requiere estadísticas de producción (ONUDI, FAO, estadísticas energéticas y de servicios) que se publican con retraso y no cubren todos los países e industrias. Cuando falta la producción, la base omite la fila doméstica. La parte de las importaciones mundiales situada en celdas país-industria con comercio doméstico observado cae del 54,9\,\% en 2015 al 45,2\,\% en 2019, el 39,9\,\% en 2021 y el 1,4\,\% en 2022. Además, trece industrias (por ejemplo, animales vivos, algodón, pesca y algunos servicios residuales) carecen de comercio doméstico en todos los países y años.

Sin corrección, el denominador de una celda afectada se reduce a las importaciones. La medida refleja entonces la concentración de la cesta importadora y no la dependencia del suministro total: un gran productor que importa poco aparecería como plenamente dependiente de su principal proveedor (por ejemplo, el crudo en Arabia Saudí o el vino en Italia). La cobertura decreciente introduciría además una tendencia artificial al alza en la dependencia. El procedimiento que sigue completa el comercio doméstico sin alterar la definición del índice.

\subsection{Estrategia de compleción}
La compleción sigue cinco principios: (i) se conserva el comercio doméstico observado, con la estrecha excepción de los residuos poco creíbles descritos en la \cref{subsec:credibilidad}; (ii) los valores ausentes se completan con el mismo concepto, la producción doméstica neta de exportaciones, tomada de las tablas input-output internacionales de la OCDE (ICIO) \citep{OECDICIO}, construidas a partir de la contabilidad nacional y complementadas con estadísticas energéticas para los productores de petróleo que no figuran en ellas; (iii) se preservan las 170 industrias y todos los países de la ITPD-E; (iv) solo se imputa donde el país importa en la industria, ya que en otro caso la diagonal no interviene en ningún cálculo; y (v) cada regla se eligió mediante validación fuera de muestra, ocultando valores observados y midiendo el error cometido al reconstruirlos.

\paragraph{Correspondencia sectorial y bloques.} La producción doméstica ICIO del país $c$ en el sector $s$ es la suma de la fila $(c,s)$ sobre las columnas del propio país $c$ (usos intermedios y finales), que equivale a la producción menos las exportaciones. Las 170 industrias de la ITPD-E se asignan a los 50 sectores ICIO (CIIU Rev.~4), tras corregir 22 asignaciones heredadas de la clasificación CIIU Rev.~3 de la ITPD-E (por ejemplo, la edición pasa a J58--J60, y el caucho y los plásticos, a C22). Las industrias que comparten sector ICIO forman un \emph{bloque} $B$ (40 en total), de modo que la suma de su comercio doméstico puede compararse con un único total ICIO, $Q_{cBt}$.

\paragraph{Calibración.} En los país-bloque-año con todas las industrias observadas, el cociente entre el comercio doméstico agregado de la ITPD-E y la producción doméstica ICIO tiene una mediana de 0,80 en manufacturas (rango intercuartílico 0,57--0,99), 1,03 en minería y energía y 1,06 en servicios. Ambas fuentes miden el mismo concepto; en manufacturas la ITPD-E es algo menor, porque las estadísticas industriales no cubren todas las empresas. El cociente mediano por bloque, $\kappa_B$, se utiliza en los controles y en los topes de bloque.

\subsection{Reglas de imputación}
Sea $D_{cit}$ el comercio doméstico del país $c$ en la industria $i$ (bloque $B$) y el año $t$, $M_{cit}$ las importaciones, $X_{cit}$ las exportaciones y $\rho_{cit} = D_{cit}/Q_{cBt}$. Los valores ausentes se completan según la jerarquía del \cref{tab:diag-reglas}.

\begin{table}[!t]
\caption{Reglas de imputación del comercio doméstico ausente}\label{tab:diag-reglas}
\centering
\footnotesize
\begin{tabularx}{\textwidth}{@{}p{1.6cm}p{4.4cm}X@{}}
\toprule
\textbf{Regla} & \textbf{Se aplica cuando} & \textbf{Estimación} \\
\midrule
a & El país tiene valores observados creíbles de la industria en un radio de 7 años (anteriores o posteriores) & $\hat D_{cit} = \bar\rho_{ci}\, Q_{cBt}$, donde $\bar\rho_{ci}$ promedia los cocientes de hasta tres años observados más próximos, ponderados por el inverso de la distancia. Los países fuera de la ICIO usan el agregado «resto del mundo» de la ICIO (a-ROW). \\
b & País ICIO sin historia propia & $\hat D_{cit} = \tilde\rho_{it}\, Q_{cBt}$, donde $\tilde\rho_{it}$ es el cociente mediano de los demás países ICIO en $t-1$, $t$ y $t+1$ (al menos tres países distintos; si no, la mediana 2016--2022). \\
b-export & Industria sin comercio doméstico observado en ningún país & $\hat D_{cit} = \kappa_B\, Q_{cBt}\, \bar X_{ci}/\bar X_{cB}$, con el peso de la industria en las exportaciones del país dentro del bloque (2016--2022) si es de al menos el 1\,\%; si no, la estructura exportadora mundial del bloque. Las industrias que forman un bloque por sí solas (p.~ej., la pesca) toman $\hat D_{cit} = Q_{cBt}$. \\
e & País fuera de la ICIO sin historia, o sin referencia para la regla b & Modelo logit de la cuota doméstica $\sigma = D/(D+M)$ con efectos fijos de industria y el logaritmo del comercio de la celda y del país (medias de tres años), más el efecto propio del país cuando tiene al menos cinco celdas observadas (e+país). $\hat D = \hat\sigma/(1-\hat\sigma)\,\bar M$, con $\hat\sigma$ acotada por los percentiles 1--99 observados en la industria. \\
en & Extracción de crudo y gas o refino de petróleo, país fuera de la ICIO con datos físicos & Producción según el \textit{Statistical Review of World Energy} \citep{EIStatReview}: crudo procesado en refinerías, o producción de crudo y gas natural en barriles equivalentes de petróleo, valorada al valor unitario mediano de los países ICIO (producción ICIO del sector por barril, escalada a la ITPD-E), por la cuota doméstica estructural de la producción del país (mediana 2016--2022 de la producción neta de exportaciones). Tiene prioridad sobre a-ROW y e. \\
\bottomrule
\end{tabularx}
\end{table}

\paragraph{Energía fuera de la ICIO.} El agregado «resto del mundo» de la ICIO mezcla economías muy distintas y orienta mal en el caso de los productores de petróleo, cuyo comercio doméstico observado en la ITPD-E es a menudo un residuo bajo o volátil. Para el crudo y el gas y para el refino de los países fuera de la ICIO, la producción se toma por ello de estadísticas energéticas sectoriales, que cubren a todos los grandes productores durante varias décadas. En una prueba que excluye cada país por turno, la regla arroja un error mediano de 4,6 puntos en la cuota doméstica del refino (5,6 en crudo y gas). Afecta en torno al 0,1\,\% de las importaciones mundiales, pero corrige resultados inverosímiles en grandes productores (por ejemplo, la cuota doméstica de los productos refinados sube del 14\,\% al 82\,\% en Irak y del 7\,\% al 39\,\% en Omán en 2022). Los valores observados por debajo del 25\,\% de la referencia energética, más de dos errores típicos por debajo de ella, se tratan como residuos.

\subsection{Credibilidad del comercio doméstico observado}\label{subsec:credibilidad}
Como el comercio doméstico es un residuo (producción menos exportaciones), deja de ser fiable cuando la producción declarada apenas supera a las exportaciones: la diferencia queda entonces dominada por el error de medida. Son ejemplos los productos farmacéuticos en Italia, que caen de 8.400 millones de dólares en 2005 a 300 millones en 2015 mientras crecen las exportaciones, o la aeronáutica en Estados Unidos, con 600 millones en 2019 entre 26.700 millones en 2018 y 52.800 millones en 2020. Usados como anclas, estos valores trasladarían cuotas domésticas casi nulas a grandes productores.

Cada celda se compara con una referencia exportadora: la producción doméstica ICIO del bloque repartida según el peso de la industria en las exportaciones del país dentro del bloque,
\begin{equation}
R_{cit} = Q_{cBt}\,\frac{X_{cit}}{X_{cBt}}.
\end{equation}
Ambas cifras de exportación proceden de la ITPD-E, por lo que las reexportaciones se compensan. Un valor observado con $D_{cit} < \varphi\, R_{cit}$ se trata como laguna y se imputa; no se usa como ancla ni en los cocientes medianos. Simétricamente, ningún valor imputado puede quedar por debajo de $\varphi\, R_{ci}$ (calculado con el peso exportador 2016--2022, para evitar saltos entre años). Se fija $\varphi = 0{,}10$, próximo al percentil 10 de $D/R$ entre las celdas observadas de grandes exportadores (mediana 0,70). El criterio afecta al 4,6\,\% de los valores domésticos observados en 2009--2022 (entre el 1,7\,\% y el 2,5\,\% de las importaciones mundiales cada año) y no altera el error fuera de muestra de las reglas (\cref{tab:diag-validacion}).

\subsection{Ajustes de coherencia}
Tras la imputación, cinco ajustes garantizan la coherencia entre años y con los totales ICIO; cada celda conserva su valor y su regla previos para asegurar la trazabilidad:
\begin{enumerate}[leftmargin=*,itemsep=2pt]
    \item \textbf{Salvaguarda de la regla a.} Si la regla a aleja la cuota doméstica más de 30 puntos porcentuales del valor observado más próximo, o da un valor más de cinco veces mayor o menor, toda la serie país-industria adopta el nivel de la observación más próxima.
    \item \textbf{Empalme entre reglas.} Cuando una serie alterna años con y sin historia propia, los segundos se reescalan para enlazar con el año más próximo basado en historia propia, conservando su propia dinámica.
    \item \textbf{Tope de magnitud.} Ningún valor de la regla b supera 1,5 veces el mayor valor observado de la industria en cualquier país (cuando hay al menos diez países observados), salvo para los tres mayores exportadores mundiales.
    \item \textbf{Topes de bloque.} Ninguna industria imputada supera $1{,}5\,\kappa_B\, Q_{cBt}$; si la suma de valores observados e imputados del bloque supera $2\,\kappa_B\, Q_{cBt}$, la parte imputada sin historia propia se reduce con un factor constante por país-bloque (mínimo del 25\,\%).
    \item \textbf{Suelo de coherencia exportadora.} Ningún valor imputado queda por debajo de $0{,}10\,R_{ci}$.
\end{enumerate}

\subsection{Validación}
El \cref{tab:diag-validacion} recoge el error fuera de muestra de las reglas seleccionadas. Se ocultan los valores domésticos observados de 2016--2021, se reconstruyen y el error se mide sobre la cuota doméstica del suministro, $\sigma = D/(D+M)$, ponderada por importaciones, que es el componente del índice afectado por la diagonal. Las reglas basadas en medianas entre países excluyen al país que se predice, y el modelo logit se valida con grupos de países no usados en la estimación.

\begin{table}[!t]
\caption{Validación fuera de muestra de las reglas de imputación (2016--2021)}\label{tab:diag-validacion}
\centering
\footnotesize
\begin{tabularx}{\textwidth}{@{}Xrr@{}}
\toprule
\textbf{Regla} & \textbf{Celdas} & \textbf{Error en $\sigma$ (pp)} \\
\midrule
a: historia propia, tres anclas, inverso de la distancia & 23.081 & 2,39 \\
\quad con anclas creíbles y suelo de coherencia exportadora & 22.303 & 2,42 \\
\quad arrastrando el último nivel (referencia) & 23.081 & 3,18 \\
a-ROW: historia propia, dinámica del resto del mundo & 10.316 & 4,31 \\
b: cociente mediano de los demás países & 23.642 & 9,65 \\
\quad agricultura / manufacturas / minería y energía / servicios & & 16,1 / 12,4 / 11,2 / 2,3 \\
b-export: estructura exportadora del país (industrias sin ninguna observación) & 23.371 & 12,37 \\
\quad estructura exportadora mundial (descartada) & 20.166 & 22,09 \\
e+país: modelo logit con efecto país (fuera de la ICIO) & 9.583 & 14,07 \\
Cuota doméstica mediana de la industria (descartada) & 9.702 & 29,61 \\
\bottomrule
\end{tabularx}
\par\smallskip\parbox{\linewidth}{\footnotesize Error: error absoluto medio en la cuota doméstica del suministro, ponderado por importaciones, en puntos porcentuales.}
\end{table}

La historia propia sigue siendo preferible a la regla entre países incluso cuando la observación más próxima tiene seis o siete años (errores de 5,0 y 8,7 puntos, frente a 10,3 y 16,1 de la regla b), lo que respalda la ventana de siete años.

\subsection{Controles de calidad}
Cada celda imputada supera siete controles con umbrales fijados de antemano: validez; coherencia de la suma del bloque con $\kappa_B\,Q_{cBt}$; cociente frente a la observación propia más próxima; continuidad interanual; cambio de la cuota doméstica respecto a la observación más próxima; verosimilitud de la cuota doméstica frente a los percentiles 1--99 observados en la industria; y magnitud respecto al mayor valor observado. Una desviación solo cuenta si el valor de referencia alternativo movería la cuota doméstica al menos cinco puntos porcentuales. Entre el 0,4\,\% y el 0,9\,\% de las importaciones imputadas quedan en la categoría roja cada año; todas esas celdas tienen una decisión documentada (por ejemplo, el mayor productor mundial de la industria o un cambio real de la producción ICIO) y ninguna queda pendiente.

\subsection{Integración en el índice}
El comercio doméstico se incorpora en un único punto, justo después de leer la ITPD-E y antes de construir las matrices; en la versión UE, antes de agregar los 27 Estados miembros, de modo que el suministro doméstico de la UE incluye el comercio doméstico de cada miembro y el comercio intra-UE. Solo se añaden filas diagonales para países e industrias ya presentes en los datos, de modo que el conjunto de países e industrias no cambia; los residuos poco creíbles se sustituyen en su lugar. El umbral de relevancia $\theta$ del cálculo (0,5\,\% del total de la columna) no se modifica; con el suministro doméstico completo, ese total incluye la producción doméstica, como ya ocurría en las celdas observadas.

\subsection{Actualización y revisiones}\label{subsec:revisiones}
El valor doméstico de un año se apoya en los años vecinos: la regla a se ancla en las observaciones más próximas, anteriores o posteriores; el cociente de la regla b usa una ventana de tres años; y algunos parámetros estructurales (los pesos exportadores, la cuota doméstica energética y $\kappa_B$) se estiman sobre una ventana de calibración, 2016--2022. Para que la serie publicada sea estable, la ventana de calibración forma parte de la versión de la metodología y no se reestima al añadir años. Cuando se dispone de los datos del año $t$, el índice se calcula para $t$ y se recalcula para $t-1$, que deja de ser el extremo de la serie; los años anteriores conservan los valores publicados. Como ocurre con el PIB y otras estadísticas económicas, las cifras de los dos últimos años son, por tanto, provisionales. Una simulación que añade 2022 a una serie que terminaba en 2021 respalda esta regla: con los parámetros fijos, cambia en torno al 1\,\% de las importaciones mundiales en 2016--2020, con un efecto despreciable en la cuota doméstica ponderada por importaciones (como máximo 0,02 puntos), mientras que en 2021 cambia el 4\,\%. La serie completa solo se recalcula y publica como nueva versión cuando cambia la metodología o cuando una nueva edición de la ITPD-E o de la ICIO revisa años pasados. Un año entra en el índice solo cuando la ICIO lo cubre.

\subsection{Limitaciones}
La regla b, la menos precisa, se aplica a cerca de un tercio de las importaciones mundiales, con errores de 16 puntos en agricultura y 12 en manufacturas, pero solo 2 en servicios. Las industrias sin ningún comercio doméstico observado dependen de supuestos estructurales (error de 12 puntos, y de hasta 27 en exportadores marginales). Los países que ensamblan exportaciones con insumos importados de la misma industria (comercio de procesamiento) muestran una dependencia alta de sus proveedores aunque sean grandes exportadores; es un rasgo de los datos de comercio por industria, no de la imputación. Por último, los últimos años de la ICIO incluyen estimaciones de la OCDE, cuyo efecto queda acotado por la salvaguarda de la regla a.

\section{Análisis de Centralidad: Índice de Hub Global (HG) / Global Hub Score (GHS)}

La identificación de intermediarios críticos es una de las mayores contribuciones del IDEE. Los países que actúan como "conectores" en las redes comerciales pueden amplificar las crisis. Este indicador de centralidad estructural se denomina indistintamente \textbf{Índice de Hub Global (HG)} en la versión en español del portal web o \textbf{Global Hub Score (GHS)} en su versión en inglés.

\subsection{Métricas de Intermediación}
Para cada país $k$, definimos dos métricas:
\begin{itemize}
    \item \textbf{Frecuencia de Intermediación ($\phi_k$):} El número de rutas comerciales significativas en las que el país $k$ participa como intermediario.
    \item \textbf{Fuerza de Intermediación ($\psi_k$):} Suma ponderada de la fuerza de dichas rutas, asignando mayor importancia a los intermediarios más cercanos al importador final.
\end{itemize}

\subsection{Índice de Centralidad Compuesto (GHS)}
El \textit{Global Hub Score} ($C_k$) integra ambas dimensiones:
\begin{equation}
C_k = \alpha \cdot \frac{\phi_k}{\max \phi} + (1-\alpha) \cdot \frac{\psi_k}{\max \psi}
\end{equation}
Empíricamente, $\alpha = 0.4$ ofrece un balance robusto, permitiendo clasificar a los países según su poder estructural como "hubs" globales que sostienen la arquitectura comercial mundial.

\section{Simulación de Disrupciones y Análisis de Resiliencia}



Para evaluar holísticamente la vulnerabilidad, implementamos simulaciones de escenarios contrafácticos. Esto implica la eliminación hipotética de nodos clave (como el suministro de un sector crítico en un país dominante) para cuantificar el impacto estructural en el resto de la red. Medimos:
\begin{itemize}
    \item \textbf{Impacto en la conectividad:} Grado en que otros países quedan aislados o sufren mermas críticas de suministro.
    \item \textbf{Reconfiguración de la centralidad:} Cómo se desplaza el poder estructural hacia otros nodos tras una disrupción.
    \item \textbf{Vulnerabilidades ocultas:} Identificación de dependencias que solo se manifiestan tras la caída de un intermediario previo.
\end{itemize}

\section{Conclusiones e Implicaciones para la Política Económica}

La principal contribución del Índice de Dependencia Económica Elcano (IDEE) es la visibilización de las vulnerabilidades "ocultas". Al descomponer la dependencia en sus componentes directos e indirectos, se facilita el diseño de respuestas políticas equilibradas.

Las dependencias directas pueden mitigarse mediante la diversificación tradicional de proveedores inmediatos. Sin embargo, las dependencias indirectas requieren una estrategia de "autonomía estratégica abierta" que incluya el fomento de capacidades domésticas en eslabones críticos del suministro y la formación de alianzas estratégicas con socios que presenten perfiles de dependencia complementarios.

Este indicador proporciona una base cuantitativa sólida para la "vigilancia estructural", permitiendo anticipar riesgos sistémicos antes de que se conviertan en crisis comerciales abiertas, y guiando el diseño de cadenas de suministro más resilientes en un entorno geopolítico incierto.

\bibliographystyle{apalike}
\bibliography{references.bib}

\end{document}